\chapter{Conclusions and Future Work}
\label{cap:conclusions}

This chapter will synthesize the findings on dataset preparation,
segmentation, geometric reference extraction, and robot operation.
It is currently a working outline: the completion notes identify the
material to add after experimental evaluation and do not report tests
or results already obtained. Detailed protocols, numerical tables, and
individual trial results belong in the experimental chapter; the
conclusions should interpret that evidence without repeating it in full.

\section{Thesis objectives and main contributions}
\label{sec:c6_objectives}

\emph{To complete:} Return to the objectives stated in the introduction
and summarize the work actually completed. Distinguish implementation
contributions from experimentally demonstrated benefits.
\begin{itemize}
    \item Dataset preparation and the use of manual and SAM~3-generated
    annotations for agricultural segmentation.
    \item Selection of a segmentation architecture and of the trained
    configuration used for robotic integration.
    \item Implementation and comparison of the two edge-based geometric
    extractors described in Chapter~\ref{cap:geometric_reference}.
    \item Integration of perception, reference extraction, and control,
    limited to the functionality and validation actually completed.
\end{itemize}

\section{Conclusions on dataset preparation and automatic annotation}
\label{sec:c6_annotation}

\emph{To complete:} Explain what the experiments establish about replacing
or reducing manual annotation. Keep direct mask-quality evaluation
separate from the performance of networks trained on those masks.
\begin{itemize}
    \item Summarize the effects of manual and automatic supervision on
    path segmentation, distinguishing the full evaluation set from the
    selected operational subset.
    \item Discuss the role of label sets and prompt targets, including
    visible soil versus a central navigation-path annotation.
    \item State whether annotation-time savings were measured, including
    inspection and correction; do not infer savings from model IoU.
    \item Identify the conditions in which automatic labels were useful
    and the cases requiring manual inspection or revision.
\end{itemize}

\section{Selection of the segmentation model for the robot}
\label{sec:c6_segmentation_selection}

\subsection{Architecture selection: accuracy and inference time}
\label{subsec:c6_architecture_selection}

\emph{To complete:} Synthesize the initial FCN, DeepLabV3, and LR-ASPP
comparison on manually annotated lettuce. Explain the practical choice
of LR-ASPP using measured accuracy and inference time under documented,
comparable conditions. Distinguish training duration from inference
latency, and model-only inference from end-to-end robot processing.

\subsection{Supervision and target-class selection}
\label{subsec:c6_supervision_selection}

\emph{To complete:} Explain the separate decision to use the binary
\texttt{allPlantsSAM} configuration. Relate its path-focused output to
the geometric extractor and summarize the evidence supporting its use
in the selected operating conditions. Do not describe it as the most
accurate configuration if the manually supervised model performs better.
State whether evaluation results also contributed to this selection.

\section{Conclusions on geometric reference extraction}
\label{sec:c6_geometry}

\emph{To complete:} Compare Method~3 and Method~4 using the evidence from
the geometric experiments, rather than their design rules alone.
\begin{itemize}
    \item Discuss interval association, independent boundary fitting,
    and centerline accuracy on the same evaluation masks.
    \item Summarize the effects of seed selection, association gaps,
    border-point exclusion, and post-fit validation.
    \item Report the implications of strip count, ROI, and parameter
    choices, distinguishing controlled comparisons from default settings.
    \item Discuss rejection behavior and incorrect but geometrically
    consistent references; a fitted line is not verified free space.
\end{itemize}

\section{Conclusions from robot field validation}
\label{sec:c6_field_validation}

\emph{To complete after the field campaign:} Summarize what closed-loop
trials establish beyond offline image segmentation and line extraction.
If field validation is not completed, explicitly state that limitation
instead of writing this section as if the trials had taken place.

\subsection{Test conditions and scope of the evidence}
\label{subsec:c6_field_conditions}

\emph{To complete:} Briefly identify the robot, crop and row conditions,
camera configuration, deployed checkpoint, geometric method, controller,
and number and length of trials. Refer to the experimental chapter for
full configuration and trial records. Distinguish any planned scenarios
from those actually tested, including illumination changes, discontinuous
rows, partial occlusions, and different initial offsets.

\subsection{Row following and task completion}
\label{subsec:c6_field_tracking}

\emph{To complete:} Interpret path-following accuracy and the fraction
of trials completed under a predefined success criterion. Summarize
travelled distance, failures, and operator interventions as well as
successful runs. Define how lateral and heading errors were measured;
image-coordinate deviations must not be presented as metric cross-track
errors without an independent measurement or justified conversion.

\subsection{Reference loss, rejection, and robot response}
\label{subsec:c6_field_failures}

\emph{To complete:} Discuss the observed behavior when the extractor
returns \texttt{degraded} or \texttt{lost}, when the reference becomes
incorrect, or when observations are interrupted. Relate these states to
the controller's implemented slowing, stopping, or recovery behavior and
to any manual intervention. Describe the supervision and protective
measures actually used in the tests. Do not infer collision avoidance
from a confidence score or from the presence of a centerline.

\subsection{Processing latency and closed-loop operation}
\label{subsec:c6_field_timing}

\emph{To complete:} Relate measured segmentation and geometric-processing
latency to the actual control update rate and robot speed. Include
end-to-end delay, variability, and dropped or stale observations where
recorded. A short network inference time alone does not establish that
the integrated system meets its timing requirements.

\section{Limitations and validity of the findings}
\label{sec:c6_limitations}

\emph{To complete:} Identify the limits supported by the final experiment
records. Resolve outstanding checks before treating them as verified facts.
\begin{itemize}
    \item Evaluation-set provenance, including the distinction between
    validation exports, full-test results, and the operational subset.
    \item Potential overlap, duplicate images, or related acquisitions
    between training and evaluation datasets.
    \item Any use of test results for model selection and the resulting
    limits on independent final evaluation.
    \item Sensitivity to label definitions, segmentation errors, image
    resolution, camera pose, and extractor parameters.
    \item Locally straight image-space boundaries, extrapolation, and
    the absence of verified obstacle detection or terrain reconstruction.
    \item Field-trial scope, measurement uncertainty, and the distinction
    between offline accuracy and closed-loop navigation performance.
\end{itemize}

\section{Future work}
\label{sec:c6_future_work}

\subsection{Dataset expansion and annotation verification}
\label{subsec:c6_future_data}

\emph{To complete:} Prioritize additional crop and lighting conditions,
manual review of difficult automatic masks, annotation-time measurement,
and acquisition-separated evaluation sets. Explain which observed
limitation motivates each proposed extension.

\subsection{Geometric extraction and temporal integration}
\label{subsec:c6_future_geometry}

\emph{To complete:} Discuss extensions justified by the measured failure
cases, such as parameter scaling, adaptive strips, temporal association,
curved references, or improved handling of partially visible boundaries.
Clearly distinguish these proposals from implemented functionality.

\subsection{Further robot trials and operational coverage}
\label{subsec:c6_future_field}

\emph{To complete:} Identify the field tests still needed to extend the
evidence: repeated trials across rows and acquisition sessions, longer
runs, additional initial offsets, and conditions not represented in the
completed campaign. Where relevant, propose evaluating reference-loss
handling and controller behavior under supervised, controlled conditions.
Do not claim operational coverage beyond the environments actually tested.

\section{Final remarks}
\label{sec:c6_final_remarks}

\emph{To complete last:} Provide a concise answer to the thesis objectives
using only findings established by the final evidence. State the level
of validation reached and the main remaining limitation. Avoid introducing
new experiments, unsupported performance claims, or results not discussed
in the preceding chapters.